\scanpage{5}
\uline{$\unsure{p|\varphi}$-structures en général.}

1) Soit $\underline S$ un site, $G$ un faisceau en groupes sur $\underline S$,
\[
f:G\mathbin{\text{---}}G
\]
un endomorphisme. On s'intéresse à la structure
suivante (où, pour $\unsure{\forall}$ torseur $P$ sous $G$,
$P^f$ désigne le torseur déduit par extension du
groupe par $f$) : donnée d'un torseur $P$
sous $G$, et d'un \uline{isom} $P\sim P^f$, ou ce qui
revient au $=$, d'un hom de faisceaux
\[
\varphi:P\longrightarrow P
\]
satisfaisant
\[
\varphi(xg)=\varphi(x)f(g).
\]
Ces objets (au dessus d'objets variables de $S$)
forment un champ\inserted{$\unsure{\Sigma}$ en groupoïdes, avec une section marquée $(G,f)$}.
On se donne des conditions
moyennant lesquelles ce dernier est une gerbe.
Noter que \uline{loc.} un tel $P$ est trivial comme
simple torseur, donc peut s'identifier à $G$ près
à une translation. La structure $\varphi$ est alors
déterminée par la connaissance de $\varphi(1)=a$,
par la formule
\[
\varphi_a(g)=a f(g)
\]
\mspage{4}
